\documentclass[12pt,a4paper]{article}

% --- PAQUETES DE IDIOMA Y FORMATO ---
\usepackage[utf8]{inputenc}
\usepackage[spanish,es-tabla]{babel}
\usepackage{geometry}
\geometry{top=2.5cm, bottom=2.5cm, left=3cm, right=2.5cm}
\usepackage{graphicx}
\usepackage{booktabs}
\usepackage{tabularx}
\usepackage{array}
\usepackage{colortbl}
\usepackage{xcolor}
\usepackage{hyperref}
\usepackage{enumitem}
\usepackage{fancyhdr}
\usepackage{microtype}
\usepackage{tcolorbox}
\usepackage{cite}

% --- DEFINICIÓN DE COLORES INSTITUCIONALES UNSAAC ---
\definecolor{unsaacBlue}{RGB}{0, 51, 102}
\definecolor{unsaacGold}{RGB}{179, 134, 0}
\definecolor{tableHeader}{RGB}{235, 242, 250}
\definecolor{tableStripe}{RGB}{248, 250, 253}
\definecolor{boxGreen}{RGB}{40, 167, 69}
\definecolor{boxBlue}{RGB}{0, 102, 179}
\definecolor{boxOrange}{RGB}{230, 126, 34}
\definecolor{boxPurple}{RGB}{142, 68, 173}

% --- CONFIGURACIÓN DE HIPERVÍNCULOS ---
\hypersetup{
    colorlinks=true,
    linkcolor=unsaacBlue,
    citecolor=unsaacBlue,
    urlcolor=unsaacBlue,
    pdfauthor={Grupo 4 - UNSAAC},
    pdftitle={Proceso de Participación en la Incubadora de Empresas Paqarina Wasi - UNSAAC}
}

% --- ENCABEZADOS Y PIES DE PÁGINA ---
\pagestyle{fancy}
\fancyhf{}
\fancyhead[L]{\small\color{gray}UNSAAC -- Ingeniería de Software I}
\fancyhead[R]{\small\color{gray}Incubadora Paqarina Wasi}
\fancyfoot[C]{\thepage}
\renewcommand{\headrulewidth}{0.4pt}

% ==============================================================================
% COMIENZO DEL DOCUMENTO
% ==============================================================================
\begin{document}

% --- PORTADA ---
\begin{titlepage}
    \centering
    \vfill
    {\large\bfseries UNIVERSIDAD NACIONAL DE SAN ANTONIO ABAD DEL CUSCO} \\[0.4cm]
    {\large FACULTAD DE INGENIERÍA ELÉCTRICA, ELECTRÓNICA, INFORMÁTICA Y MECÁNICA} \\[0.4cm]
    {\large ESCUELA PROFESIONAL DE INGENIERÍA INFORMÁTICA Y DE SISTEMAS} \\[1.5cm]
    
    \includegraphics[width=0.35\textwidth]{unsaac.jpg} \\[1.5cm] % Nota: Cambiar por logo local si es necesario, o Overleaf descargará automáticamente.
    
    {\Large\bfseries ASIGNATURA: INGENIERÍA DE SOFTWARE I} \\[0.6cm]
    {\large\bfseries TRABAJO: PROCESO INTEGRAL DE PARTICIPACIÓN Y PROGRAMAS DE EMPRENDIMIENTO EN LA INCUBADORA DE EMPRESAS PAQARINA WASI - UNSAAC} \\[0.6cm]
    {\large\bfseries RUTA DE ACOMPAÑAMIENTO: SENSIBILIZACIÓN, IDEATHON, HACKATHON, PREINCUBACIÓN E INCUBACIÓN} \\[1.5cm]
    
    \textbf{DOCENTE RESPONSABLE:} \\
    Mgt. Luz Indira Ticona Felix \\[1.0cm]
    
    \textbf{INTEGRANTES / DESARROLLADORES (GRUPO 4):} \\
    ALVAREZ CATUNTA, Angel Ismael - 221945 \\
    CHILLIHUANI HUAMAN, Rivaldo Franco - 221445 \\
    SUPA CUSIPAUCAR, Yeferson - 220553 \\
    VILLAVICENCIO SEGUIL, Edu Piero - 220555 \\[0.5cm]
    
    {\large Cusco - Perú} \\
    {\large 2026}
    \vfill
\end{titlepage}

\newpage
\tableofcontents
\newpage

% ==============================================================================
% 1. INTRODUCCIÓN Y CONTEXTO INSTITUCIONAL
% ==============================================================================
\section{Introducción y Contexto Institucional}

\subsection{¿Qué es Paqarina Wasi?}
\textbf{Paqarina Wasi} (expresión quechua que significa \textit{``Casa del Renacer''} o \textit{``Casa del Nacimiento''}) es la \textbf{Incubadora de Empresas de Base Tecnológica e Innovación} de la \textbf{Universidad Nacional de San Antonio Abad del Cusco (UNSAAC)}. Opera bajo la dirección del Vicerrectorado de Investigación (VRIN) y cuenta con el respaldo estratégico y cofinanciamiento del programa nacional \textbf{ProInnóvate} del Ministerio de la Producción (PRODUCE) del Perú.

Su misión principal es dinamizar el ecosistema de emprendimiento e innovación en la región del Cusco, transformando el conocimiento científico, las investigaciones formativas y los desarrollos tecnológicos generados en las aulas y laboratorios universitarios en empresas sostenibles, escalables y con alto impacto económico y social (\textit{startups}).

\subsection{Importancia para los Estudiantes de Ingeniería Informática y de Sistemas}
En la formación profesional de los estudiantes de Ingeniería Informática y de Sistemas, el desarrollo de software no concluye en la codificación de un producto mínimo; requiere un modelo de negocio sostenible, validación en el mercado real y mecanismos de financiamiento. Paqarina Wasi actúa como el puente indispensable para que proyectos de software académico (tales como sistemas de optimización de recolección de residuos, plataformas IoT, soluciones de inteligencia artificial o aplicaciones móviles) transiten desde un prototipo de laboratorio hacia una empresa formal de base tecnológica.

% ==============================================================================
% 2. PROGRAMAS Y RUTA DE ACOMPAÑAMIENTO
% ==============================================================================
\section{Programas y Rutas de Acompañamiento en Paqarina Wasi}

Paqarina Wasi estructura su propuesta de valor a través de una ruta progresiva (embudo de emprendimiento), diseñada para atender a los participantes según la madurez de su proyecto:

\subsection{1. Sensibilización (Etapa de Inspiración y Cultura Emprendedora)}
\begin{itemize}[leftmargin=*]
    \item \textbf{Propósito:} Despertar el interés, la curiosidad y la mentalidad emprendedora en la comunidad universitaria y el público en general.
    \item \textbf{Duración:} 2 a 4 horas.
    \item \textbf{Perfil del Participante:} Estudiantes, docentes y profesionales sin experiencia previa en negocios.
    \item \textbf{Dinámica:} Conferencias magistrales, testimonios de fundadores de startups peruanas, dinámicas de innovación abierta y espacios de networking inicial.
\end{itemize}

\subsection{2. Ideathon (Etapa de Creatividad Multidisciplinaria)}
\begin{itemize}[leftmargin=*]
    \item \textbf{Propósito:} Generar ideas innovadoras frente a desafíos regionales específicos (turismo, agroindustria, transporte, medio ambiente).
    \item \textbf{Duración:} 1 jornada intensiva (1 día).
    \item \textbf{Perfil del Participante:} Estudiantes y profesionales de todas las facultades. No se requieren conocimientos técnicos de programación.
    \item \textbf{Metodología:} Talleres rápidos de \textit{Design Thinking}, lluvias de ideas estructuradas, validación cualitativa y presentación final en formato \textit{pitch} relámpago.
\end{itemize}

\subsection{3. Hackathon (Etapa de Creación Rápida de Prototipos / MVP)}
\begin{itemize}[leftmargin=*]
    \item \textbf{Propósito:} Construir soluciones tecnológicas funcionales (software, hardware o IoT) para resolver problemáticas complejas en un lapso reducido.
    \item \textbf{Duración:} 48 a 72 horas ininterrumpidas de desarrollo.
    \item \textbf{Perfil del Participante:} Equipos multidisciplinarios compuestos por programadores, diseñadores UX/UI, administradores y especialistas temáticos.
    \item \textbf{Resultados Esperados:} Validación técnica de la arquitectura de software, conformación de equipos cofundadores y un Producto Mínimo Viable (MVP) listo para postular a preincubación.
\end{itemize}

\subsection{4. Programa de Preincubación (Validación y Estructuración)}
\begin{itemize}[leftmargin=*]
    \item \textbf{Propósito:} Validar el problema y la solución en el mercado real, y estructurar formalmente el modelo de negocio.
    \item \textbf{Duración:} 2 meses.
    \item \textbf{Enfoque:} Validación temprana y ajuste problema-solución (\textit{Problem-Solution Fit}).
    \item \textbf{Beneficios Clave:}
    \begin{itemize}
        \item Asignación de mentores especializados en modelos de negocio Canvas y Lean Startup.
        \item Acceso libre a espacios de \textit{coworking} dentro del campus universitario.
        \item Talleres de prototipado rápido, entrevistas con clientes y validación cuantitativa.
        \item Puente directo y preparación para el programa de incubación formal.
    \end{itemize}
\end{itemize}

\subsection{5. Programa de Incubación (Consolidación y Escalamiento Empresarial)}
\begin{itemize}[leftmargin=*]
    \item \textbf{Propósito:} Consolidar startups con productos mínimos validados y equipos constituidos, logrando tracción comercial y levantamiento de capital.
    \item \textbf{Duración:} 2 meses intensivos.
    \item \textbf{Enfoque:} Crecimiento, penetración de mercado (\textit{Product-Market Fit}) y escalabilidad.
    \item \textbf{Beneficios Clave:}
    \begin{itemize}
        \item Mentorías individuales con especialistas sectoriales (legal, tributario, finanzas, arquitectura de software).
        \item Acompañamiento especializado para postular a fondos concursables no reembolsables de \textbf{ProInnóvate (hasta S/. 67,000} mediante StartUp Perú - Emprendimientos Innovadores).
        \item Acceso a laboratorios especializados de la UNSAAC y herramientas en la nube.
        \item Preparación para el \textbf{Demo Day}, conectando directamente con redes de inversionistas ángeles nacionales e internacionales.
    \end{itemize}
\end{itemize}

\subsection{6. Programa de Aceleración (\textit{Próximamente})}
\begin{itemize}[leftmargin=*]
    \item \textbf{Propósito:} Escalamiento exponencial e internacionalización de startups consolidadas con ventas recurrentes y tracción probada.
    \item \textbf{Duración:} 3 a 4 meses.
    \item \textbf{Enfoque:} Rondas de inversión Serie A/B, optimización de métricas CAC/LTV y apertura de mercados globales.
\end{itemize}

% ==============================================================================
% 3. PROCESO PASO A PASO PARA PARTICIPAR
% ==============================================================================
\newpage
\section{Proceso Paso a Paso para Participar en Paqarina Wasi}

Para que un equipo de estudiantes o emprendedores pueda integrarse formalmente a Paqarina Wasi, se debe seguir la siguiente ruta operativa:

\begin{tcolorbox}[colback=white,colframe=boxBlue,title={\textbf{Fase 1: Convocatoria y Descarga de Bases}},fonttitle=\bfseries]
Las convocatorias se lanzan de manera periódica (generalmente al inicio de cada semestre académico, por ejemplo bajo programas como \textit{Cusco Startup Lab} o convocatorias directas de Pre-incubación). Los postulantes deben revisar las bases oficiales publicadas en el portal institucional de la UNSAAC y en la página oficial de Facebook e Instagram (\textbf{@PaqarinaWasiUNSAAC}).
\end{tcolorbox}

\begin{tcolorbox}[colback=white,colframe=boxGreen,title={\textbf{Fase 2: Conformación del Equipo Multidisciplinario}},fonttitle=\bfseries]
Las bases estipulan que los proyectos deben estar conformados por equipos de 2 a 5 integrantes. Al menos uno o más integrantes deben pertenecer activamente a la comunidad antoniana (estudiantes de pregrado, posgrado, docentes o egresados de la UNSAAC). Se valora positivamente la complementariedad entre perfiles tecnológicos (Ingeniería Informática/Sistemas) y comerciales (Administración, Economía, etc.).
\end{tcolorbox}

\begin{tcolorbox}[colback=white,colframe=boxOrange,title={\textbf{Fase 3: Registro y Formulario de Postulación}},fonttitle=\bfseries]
El equipo completa el formulario digital de inscripción dentro de los plazos establecidos. En esta solicitud se documenta:
\begin{itemize}[leftmargin=*]
    \item \textbf{Datos del Equipo:} Roles, horas de dedicación semanal y vínculos con la universidad.
    \item \textbf{Problema Detectado:} Descripción clara del dolor de mercado o necesidad insatisfecha.
    \item \textbf{Propuesta de Solución:} Explicación técnica de la solución (software, plataforma, app o hardware) y su factor innovador/diferenciador.
    \item \textbf{Grado de Madurez:} Si se postula a Preincubación (idea con boceto o prototipo conceptual) o a Incubación (MVP operativo con pruebas iniciales).
\end{itemize}
\end{tcolorbox}

\begin{tcolorbox}[colback=white,colframe=boxPurple,title={\textbf{Fase 4: Evaluación Técnica y Filtro de Elegibilidad}},fonttitle=\bfseries]
El comité evaluador de Paqarina Wasi verifica el cumplimiento de los requisitos formales (elegibilidad de integrantes y originalidad del proyecto) y califica la propuesta según una matriz de evaluación estandarizada (grado de innovación, tamaño de mercado, viabilidad técnica y compromiso del equipo).
\end{tcolorbox}

\begin{tcolorbox}[colback=white,colframe=tierAStar,title={\textbf{Fase 5: Pitch Day y Selección Final}},fonttitle=\bfseries]
Los proyectos preseleccionados son convocados a una jornada de \textbf{Pitch Day}. Cada equipo dispone de un tiempo delimitado (generalmente de 3 a 5 minutos de exposición seguidos de 3 minutos de preguntas) ante un jurado calificador compuesto por mentores, directores de la incubadora y representantes del ecosistema de innovación. Los equipos seleccionados son anunciados oficialmente mediante acta resolutiva.
\end{tcolorbox}

\begin{tcolorbox}[colback=white,colframe=unsaacBlue,title={\textbf{Fase 6: Firma de Acta de Compromiso e Inicio}},fonttitle=\bfseries]
Los equipos admitidos firman una carta o convenio de compromiso con Paqarina Wasi donde se estipula la asistencia a los talleres, el cumplimiento de los hitos periódicos y el uso de los espacios de coworking. A partir de ese momento, inician la ruta lectiva y mentorías personalizadas.
\end{tcolorbox}

% ==============================================================================
% 4. MATRIZ RESUMEN DE LOS PROGRAMAS
% ==============================================================================
\newpage
\section{Matriz Comparativa de los Programas de Paqarina Wasi}

La siguiente tabla sintetiza las características, perfiles de ingreso y resultados esperados de cada programa ofrecido por la incubadora:

\begin{table}[htbp]
\centering
\small
\renewcommand{\arraystretch}{1.3}
\begin{tabularx}{\textwidth}{p{2.5cm} c p{2.8cm} p{3.2cm} >{\raggedright\arraybackslash}X}
\toprule
\rowcolor{tableHeader}
\textbf{Programa} & \textbf{Duración} & \textbf{Perfil de Entrada} & \textbf{Metodología Clave} & \textbf{Resultado / Entregable} \\
\midrule
\rowcolor{tableStripe}
\textbf{Sensibilización} & 2 a 4 horas & Abierto a todo público sin experiencia. & Charlas magistrales y dinámicas de grupo. & Inspiración y acercamiento al ecosistema innovador. \\
\textbf{Ideathon} & 1 día & Equipos de cualquier facultad o carrera. & \textit{Design Thinking}, mapas de empatía y pitch. & Idea estructurada y validada a nivel conceptual. \\
\rowcolor{tableStripe}
\textbf{Hackathon} & 48 a 72 h & Programadores, diseñadores y gestores. & Maratón intensiva de desarrollo de código. & Prototipo funcional (MVP) y equipo cofundador. \\
\textbf{Preincubación} & 2 meses & Emprendedores en fase de idea estructurada. & \textit{Lean Canvas}, mentorías y entrevistas. & Modelo de negocio validado con usuarios reales. \\
\rowcolor{tableStripe}
\textbf{Incubación} & 2 meses & Startups con MVP y equipo consolidado. & \textit{Go-to-market}, finanzas y rondas de inversión. & Tracción, Demo Day y postulación a ProInnóvate (S/. 67k). \\
\textbf{Aceleración} & 3 a 4 meses & Startups en ventas y tracción constante. & Escalamiento global y negociación Serie A/B. & Expansión internacional y levantamiento de capital. \\
\bottomrule
\end{tabularx}
\caption{Comparativa integral de los programas de emprendimiento en Paqarina Wasi - UNSAAC.}
\label{tab:paqarina_programas}
\end{table}

% ==============================================================================
% 5. REQUISITOS GENERALES Y CRITERIOS DE EVALUACIÓN
% ==============================================================================
\section{Requisitos Generales y Criterios de Evaluación}

\subsection{Requisitos de Elegibilidad}
\begin{enumerate}[leftmargin=*]
    \item \textbf{Vínculo Institucional:} Al menos un miembro fundador debe ser estudiante regular, egresado, docente o investigador de la UNSAAC.
    \item \textbf{Composición del Equipo:} Mínimo 2 integrantes, recomendándose un equipo multidisciplinario (por ejemplo: un líder de desarrollo de software, un diseñador de experiencia de usuario y un responsable de gestión comercial).
    \item \textbf{Disponibilidad de Tiempo:} Compromiso documentado de asistir a los talleres semanales y sesiones de mentoría especializada.
    \item \textbf{Originalidad:} La solución propuesta debe ser original o presentar un grado sustancial de diferenciación frente a soluciones existentes en el mercado.
\end{enumerate}

\subsection{Criterios de Evaluación en las Convocatorias}
\begin{itemize}[leftmargin=*]
    \item \textbf{Grado de Innovación y Diferenciación (30\%):} Se evalúa si el producto aporta un valor tecnológico novedoso frente a las soluciones comerciales tradicionales.
    \item \textbf{Potencial de Mercado y Escalabilidad (25\%):} Capacidad del modelo de negocio de crecer a nivel regional, nacional o internacional sin un incremento proporcional lineal de costos.
    \item \textbf{Capacidad del Equipo Emprendedor (25\%):} Competencias técnicas, compromiso demostrado y sinergia entre los integrantes para llevar a cabo la ejecución del proyecto.
    \item \textbf{Viabilidad Técnica y Nivel de Madurez (20\%):} Factibilidad real de desarrollo del software o prototipo dentro de los plazos establecidos.
\end{itemize}

% ==============================================================================
% 6. CANALES OFICIALES DE COMUNICACIÓN Y CONTACTO
% ==============================================================================
\section{Canales Oficiales de Atención y Monitoreo}

Para estar informado sobre las fechas de apertura de convocatorias, los postulantes deben seguir los siguientes canales institucionales:

\begin{itemize}[leftmargin=*]
    \item \textbf{Redes Sociales Oficiales:}
    \begin{itemize}
        \item Facebook: \href{https://www.facebook.com/PaqarinaWasiUNSAAC}{Paqarina Wasi - Incubadora de Empresas UNSAAC}
        \item Instagram: \texttt{@paqarinawasi\_unsaac}
    \end{itemize}
    \item \textbf{Correo Electrónico de Contacto:} \href{mailto:paqarinawasi@unsaac.edu.pe}{\texttt{paqarinawasi@unsaac.edu.pe}}
    \item \textbf{Ubicación Física:} Campus Universitario de Perayoc, Universidad Nacional de San Antonio Abad del Cusco, Pabellón de Innovación y Coworking.
    \item \textbf{Portal Institucional UNSAAC:} Sección de Convocatorias y Emprendimiento en \url{https://www.unsaac.edu.pe}.
\end{itemize}

% ==============================================================================
% 7. CONCLUSIONES Y RECOMENDACIONES
% ==============================================================================
\section{Conclusiones y Recomendaciones}

\begin{itemize}[leftmargin=*]
    \item \textbf{Ecosistema Articulado:} Paqarina Wasi brinda una estructura progresiva y lógica que permite a un estudiante de ingeniería pasar desde una idea inicial generada en un Hackathon o Ideathon hasta una startup formal con financiamiento de ProInnóvate.
    \item \textbf{Oportunidad para Proyectos Universitarios:} Trabajos desarrollados en asignaturas como Ingeniería de Software I y II (como por ejemplo aplicaciones móviles, plataformas web y sistemas de optimización urbana) cuentan con un potencial inmediato para postular a los programas de \textit{Preincubación}.
    \item \textbf{Recomendación Estratégica:} Se recomienda a los equipos estudiantiles no postular de forma individual con un perfil netamente técnico; la conformación de equipos interdisciplinarios (incorporando estudiantes de carreras como Administración, Contabilidad o Marketing) incrementa significativamente la tasa de éxito en el \textit{Pitch Day} y la posterior adjudicación de fondos estatales.
\end{itemize}

% ==============================================================================
% REFERENCIAS BIBLIOGRÁFICAS
% ==============================================================================
\newpage
\begin{thebibliography}{99}

\bibitem{proinnovate}
Ministerio de la Producción del Perú -- ProInnóvate. (2024). 
\textit{Bases del Concurso StartUp Perú: Capital Semilla para Emprendimientos Innovadores}. 
Gobierno del Perú. Disponible en: \url{https://www.proinnovate.gob.pe}

\bibitem{leyuniversitaria}
Congreso de la República del Perú. (2014). 
\textit{Ley Universitaria N° 30220 -- De los Centros de Producción, Innovación e Incubadoras de Empresas}. 
Diario Oficial El Peruano.

\bibitem{unsaac_vrin}
Vicerrectorado de Investigación (VRIN-UNSAAC). (2025). 
\textit{Reglamento General de Emprendimiento, Transferencia Tecnológica e Incubación de Empresas}. 
Universidad Nacional de San Antonio Abad del Cusco.

\bibitem{ries}
Ries, E. (2011). 
\textit{The Lean Startup: How Today's Entrepreneurs Use Continuous Innovation to Create Radically Successful Businesses}. 
Crown Business.

\bibitem{osterwalder}
Osterwalder, A., \& Pigneur, Y. (2010). 
\textit{Business Model Generation: A Handbook for Visionaries, Game Changers, and Challengers}. 
John Wiley \& Sons.

\end{thebibliography}

\end{document}
